\documentclass[11pt]{article}
\usepackage[margin=1in]{geometry}
\usepackage[T1]{fontenc}
\usepackage{lmodern}
\usepackage{amsmath}
\usepackage{graphicx}
\usepackage{subfig}
\usepackage{cleveref}

% subfig (Cohen) is the older sibling of subcaption and uses a different
% command surface: \subfloat[<caption>]{<body>} rather than the subfigure
% environment. This fixture exercises that surface together with
% \ContinuedFloat, which lets a second float reuse the first float's number.

\crefname{subfigure}{panel}{panels}
\Crefname{subfigure}{Panel}{Panels}

\title{Panelized Displays with the \texttt{subfig} Package}
\author{Test Corpus}
\date{}

\begin{document}
\maketitle

\section{Introduction}

This fixture isolates the \texttt{subfig} package, which predates
\texttt{subcaption} and remains common in legacy manuscripts and in journal
templates that have not migrated. The two packages are mutually exclusive and
expose different syntax: \texttt{subfig} provides the \verb|\subfloat| command,
whereas \texttt{subcaption} provides a \texttt{subfigure} environment. A
converter that special-cases only the environment form will silently drop the
panels produced here, so the fixture is meant to surface that gap. None of the
panels depend on external raster files; each is a fixed-size \verb|\framebox|
whose interior strut reserves vertical space, which keeps the document
self-contained under \texttt{pdflatex}.

\Cref{fig:panels} collects three panels in a single figure float. Each panel
carries an independent caption and label, while the enclosing float supplies a
common caption and its own label. The companion display in
\cref{fig:panels-cont} reuses the same figure number through
\verb|\ContinuedFloat|, so the two appear to a reader as one logical figure
split across a page boundary.

\section{A three-panel float}

\begin{figure}[htbp]
\centering
\subfloat[Increasing series\label{fig:panel-a}]{%
  \framebox[3.4cm][c]{\rule{0pt}{2.6cm}Panel A}}
\hfill
\subfloat[Grouped bars\label{fig:panel-b}]{%
  \framebox[3.4cm][c]{\rule{0pt}{2.6cm}Panel B}}
\hfill
\subfloat[Residual scatter\label{fig:panel-c}]{%
  \framebox[3.4cm][c]{\rule{0pt}{2.6cm}Panel C}}
\caption{Three placeholder panels typeset with \texttt{subfig}. The panel
  bodies are fixed-width boxes rather than images so the fixture compiles
  without any external graphics.}
\label{fig:panels}
\end{figure}

The panels are cross-referenced individually to verify that subpanel labels
resolve through \texttt{cleveref}. \Cref{fig:panel-a} shows the increasing
series, \cref{fig:panel-b} shows the grouped bars, and \cref{fig:panel-c} shows
the residual scatter. A combined reference such as
\cref{fig:panel-a,fig:panel-b,fig:panel-c} should expand into a list of panel
identifiers, confirming that the subpanel counter and the \texttt{cleveref}
name mapping agree.

\section{A continued float}

\begin{figure}[htbp]
\ContinuedFloat
\centering
\subfloat[Diagnostic histogram\label{fig:panel-d}]{%
  \framebox[3.4cm][c]{\rule{0pt}{2.6cm}Panel D}}
\hfill
\subfloat[Quantile plot\label{fig:panel-e}]{%
  \framebox[3.4cm][c]{\rule{0pt}{2.6cm}Panel E}}
\caption{Continuation of the preceding figure. Because the float begins with
  \texttt{\textbackslash ContinuedFloat}, it shares the figure number of
  \cref{fig:panels} and is conventionally read as ``(cont.)''.}
\label{fig:panels-cont}
\end{figure}

\section{Discussion}

The continuation in \cref{fig:panels-cont} adds two further panels,
\cref{fig:panel-d} and \cref{fig:panel-e}, under the shared figure number. The
substantive point for a fidelity test is that the panel-level structure and the
two-part numbering both carry meaning that a flat extraction would lose: a
reader who sees only ``Figure~1'' and ``Figure~1 (cont.)'' still expects five
labelled panels, A through E, to be addressable in the surrounding prose. When
the panels themselves are images, the same structure governs how alt text and
captions should be associated, so preserving the \verb|\subfloat| grouping
matters independently of whether the bodies are boxes or figures.

\end{document}
